\section{Degree-four tensors in spin families}
\label{sec:auxiliary}

The exact Kuga--Satake correspondences control each K3
surface separately.  To compare several factors, we will
also need a quadratic contraction tensor on certain
abelian varieties.  The construction below permits
independent period planes at every real place of a
totally real field.  Its four weight-one slots will
later compare actions on different spin constituents.

Let $F'/\Q$ be a totally real Galois extension and let
$(V',q')$ be a quadratic space over $F'$ of dimension $d'$.
Assume that
\begin{equation}
 \label{aux:quadratic-assumptions}
 d'\ge5,\quad d'\equiv5\pmod8,\qquad
 \operatorname{sig}(V'_\sigma)=(2,d'-2)
       \quad(\sigma:F'\hookrightarrow\R),\qquad
 \operatorname{Witt}(V')\ge2.
\end{equation}
The dimension can be increased as needed in the
application.  Choose an
oriented positive period plane at each real place.  The
corresponding spin circles give the weight-one Hodge
structure
\begin{equation}
 \label{aux:weight-one}
 W'=\operatorname{Res}_{F'/\Q}C^+_{F'}(V').
\end{equation}
It is polarizable: choose an invariant alternating form
with the required positivity at every real embedding,
using approximation over $F'$ in the construction
\eqref{ks:trace-form}, and take its trace to $\Q$.
Thus $W'$ is the rational first cohomology of an abelian
variety $A'$, with a chosen lattice.  On each component
the period construction gives an algebraic family of
abelian varieties.  Indeed, a neat arithmetic subgroup
of the restriction-of-scalars spin group can be chosen
to preserve the lattice and polarization.  The product
of the orthogonal period domains then maps
holomorphically and equivariantly to Siegel space, and
the algebraicity theorems used in
Section~\ref{ks:deformation} apply.

Choose a volume element $z'$ over $F'$.  Define a
rational tensor
\begin{equation}
 \label{aux:P}
 \mathcal P\in(\End W')^{\otimes2}
\end{equation}
by prescribing its extension to $\C$.  For each embedding
$\sigma:F'\hookrightarrow\C$, take dual bases
$(v_j)$ and $(v_j^\vee)$ of $V'_\sigma$ and form
\[
 \sum_jL_{v_jz'}\otimes L_{v_j^\vee z'}.
\]
Each operator is supported on the $\sigma$ summand of
$W'_\C$; sum these tensors over $\sigma$.  This is the
trace construction, so it descends to $\Q$.  It is
invariant under $\operatorname{Res}_{F'/\Q}\Spin(V')$.
Using an invariant polarization in each endomorphism
factor, regard $\mathcal P$ as a degree-four Hodge
tensor in four ordered $H^1(A')$ slots.

\begin{proposition}
 \label{prop:auxiliary-spin}
For the data \eqref{aux:quadratic-assumptions}, the tensor
$\mathcal P$ is the class of a cycle in
$\CH^2((A')^4)_{\Q}$, for every choice of the period planes and
every isogenous realization of \eqref{aux:weight-one}.
\end{proposition}

\begin{proof}
We first choose a point of the family whose abelian
variety is a power of the non-CM elliptic curve $E$
in the torus dictionary.  The Witt-index assumption
gives a split rank-three quadratic subspace of $V'$
with signature $(2,1)$ at every real place: take one
hyperbolic plane and a vector positive at all places
in a second one.  Its spin group is $\mathrm{SL}_2$
over $F'$.  Restricting scalars and then taking the
diagonal $\mathrm{SL}_2$ over $\Q$ makes $W'$ a
multiple of its standard representation.

The common circle can be placed in any prescribed
component of the period family.  Under the ternary
identification, conjugate by an element of
$\mathrm{GL}_2(F')$ with the chosen determinant signs
at the real embeddings.  Approximation permits these
signs to be chosen independently.  This also accommodates
either elliptic marking convention in the graph
dictionary; its common line sign was fixed before the
choice of tensor.  Choose the lattice to give the
corresponding elliptic power exactly at this point.

Apply the torus-only realization with
\[
 W=(W')^{\oplus4}=M_0\otimes H,
\]
using bases compatible with the four copies.  Place
$\mathcal P$ in $\bigwedge^4W$, with one vector in
each copy, after the polarization conversion.  These copies
have a second use: to verify the wedge-injectivity hypothesis,
we will project a degree-one class from one copy against an
annihilator's degree-two component on another.  We arrange
that the annihilator has the same component on every copy.

Let
$\theta$ be the direct sum of four identical copies
of a nondegenerate rational Spin-invariant alternating
tensor on $W'$.  Choose two more tensors
$\theta_1,\theta_2$ by tensoring that tensor with
nondegenerate symmetric matrices on the four copy
indices: the first is diagonal with distinct entries,
and the second has every entry nonzero.  Here the relative
operators act on the four copy indices.  Their commutation
conditions will force identical blocks there, while allowing
operators within $W'$; this is the copy condition just required.
As in the earlier construction, separate the powers so that
wedge, derivation, and contraction terms have different degrees.
Put
\begin{equation}
 \label{aux:test-tensor}
 C_* =\mathcal P+\theta^6+\theta_1^9+\theta_2^{12}.
\end{equation}
This is Spin-invariant and hence invariant under the
chosen simultaneous $\mathrm{SL}_2$.  The torus
dictionary produces a middle-degree primitive class
$a$.  Here
\[
 g=\tfrac12\dim_{\Q}W=[F':\Q]2^{d'}\ge32,
\]
so $g$ is even, the realization dimension condition is
satisfied, and the symplectic wedge maps used below
are in their injective range.

We verify the wedge-injectivity hypothesis of
Proposition~\ref{prop:realization}.  In the formal
annihilation equation $F(\lambda,L,\pi)C_*=0$, the
three operator types on the displayed powers have
degrees $(10,12,14)$, $(16,18,20)$, and $(22,24,26)$.
They are all distinct and exceed the degrees $2,4,6$
obtained by acting on $\mathcal P$.  The same symplectic
wedge and contraction calculation as in
Lemma~\ref{ks:formal-kernel} therefore gives
$\lambda=\pi=0$ and
\[
 F_L(\theta)=F_L(\theta_1)=F_L(\theta_2)=0.
\]
The first condition makes $L$ self-adjoint.  Commutation
with the diagonal relative operator makes it block
diagonal on the four copies.  Commutation with the
second relative operator, whose off-diagonal entries
are nonzero, then makes all four blocks equal.

Under the inverse dictionary, the annihilator
$\mathcal Z$ consequently consists of degree-two forms
obtained by repeating the same form on each of the
four torus copies.  Its projection to the degree-two
cohomology of any one copy is injective.  For an input
in $H^1$ of copy $i$, choose $j\ne i$ and project its
wedge product to multidegree one on copy $i$ and
degree two on copy $j$.  This is the tensor product
of the identity on $H^1$ of copy $i$ with that
injective projection of $\mathcal Z$.  Inputs from
the other copies have different multidegrees.  Taking
these projections for all $i$ proves
\eqref{setup:wedge-injective}.

In the torus-only case, the Euler tests determine
the entire normalized Mukai vector, so the resulting
perfect complex has $C=C_*$.  At this elliptic point,
the actual cap kernel and the formal annihilator
have the same dimension: the two sets of weights
on this invariant tensor are conjugate over $\C$.
The Ext bound is therefore an equality at the cap
rank, and the semiregularity trace is injective.

Throughout the indicated algebraic period family,
all terms of \eqref{aux:test-tensor} remain parallel
Hodge classes by Spin invariance.  Apply
Proposition~\ref{prop:semiregularity} and extract
cohomological degree four.  The three higher powers
have degrees $12$, $18$, and $24$, so the extracted
class is exactly $\mathcal P$.  Its one-slot-per-copy
component gives an algebraic class on $(A')^4$.
The elliptic starting point was available in every
component.  Finally, transport under abelian
isogenies and change of rational polarization use
algebraic correspondences, proving the assertion
for all the stipulated realizations.
\end{proof}
